\chapter{اللقاء السابع والعشرون: السلم كافة وأصول الاختلاف وبدايات آيات القتال}

\section{مقدمة}
السلام عليكم ورحمة الله وبركاته. الحمد لله رب العالمين، وأشهد أن لا إله إلا الله وحده لا شريك له، وأشهد أن محمداً عبده ورسوله. اللهم صل على محمد وعلى آل محمد كما صليت على آل إبراهيم إنك حميد مجيد، اللهم بارك على محمد وعلى آل محمد كما باركت على آل إبراهيم إنك حميد مجيد.

استأنف الشيخ هذا اللقاء من قوله تعالى: \begin{ayah}يَا أَيُّهَا الَّذِينَ آمَنُوا ادْخُلُوا فِي السِّلْمِ كَافَّةً\end{ayah}، ثم مضى في آيات التحذير من الزلل بعد البينات، وذكر سنن الأمم في الاختلاف، وشيئاً من أحكام القتال في الشهر الحرام، حتى انتهى إلى قوله تعالى: \begin{ayah}إِنَّ الَّذِينَ آمَنُوا وَالَّذِينَ هَاجَرُوا وَجَاهَدُوا فِي سَبِيلِ اللَّهِ\end{ayah}.

\separator

\section{تأويل (يا أيها الذين آمنوا ادخلوا في السلم كافة)}
\isnad{قال أبو جعفر:} اختلف أهل التأويل في معنى \textit{السلم} في هذا الموضع:
\begin{itemize}
    \item فقيل: هو الإسلام.
    \item وقيل: هو الطاعة.
    \item وقيل عند بعض أهل القراءة: هو الصلح والمسالمة.
\end{itemize}

ثم ذكر \rawi{الطبري} اختلاف القراء في فتح السين وكسرها، ووازن بين دلالة القراءتين، ثم رجح أن المراد: ادخلوا في الإسلام كافة، أي: التزموا شرائعه كلها، ولا تأخذوا ببعضها دون بعض.

ونبّه الشيخ هنا إلى أمر مهم في قراءة صنيع \rawi{الطبري} في القراءات: وهو أن ترجيحه بين القراءات لا يعني دائماً رد الأخرى أو إبطالها، بل قد يكون مراده بيان الأليق بالسياق والأقرب في المعنى.

\begin{fawaid}[ينبغي التفريق بين الترجيح بين القراءات ورد القراءة]
ليس كل موضع رجح فيه الطبري قراءة على أخرى يكون مراده إبطال القراءة الأخرى؛ بل كثيراً ما يكون المقصود بيان الأوجه في المعنى والسياق. والتحرير في هذا الباب يقي من إساءة الظن بالأئمة أو اختزال مناهجهم في عبارة مجملة.
\end{fawaid}

\section{وجه الخطاب في الآية}
ذكر \rawi{الطبري} احتمالين في المخاطبين بهذه الآية:
\begin{itemize}
    \item أن يكونوا المؤمنين بمحمد صلى الله عليه وسلم، فيكون المعنى أمراً لهم باستكمال شرائع الإسلام كلها.
    \item أو أن يكونوا من أهل الإيمان بالأنبياء قبل محمد صلى الله عليه وسلم ممن لم يدخلوا في الإيمان به، فيكون المعنى دعوة لهم إلى الدخول في الإسلام التام.
\end{itemize}

ثم بنى على هذا الترجيح رده للقول الذي حمل \textit{السلم} على الصلح؛ لأن السياق عنده لا يساعد عليه، ولأن نظائر القرآن لا تشهد له في هذا الموضع.

\begin{fawaid}[القول لا يرد بالهوى بل ببيان حجته]
من أنفع ما يربي عليه تفسير الطبري أن القارئ لا يكتفي بمعرفة القول المختار، بل يتعلم كيف يقيم العالم حجته على الترجيح والرد، ولو خالفناه بعد ذلك في النتيجة. ففهم مأخذ القول جزء من العلم، لا أمراً زائداً عليه.
\end{fawaid}

\separator

\section{تأويل (فإن زللتم من بعد ما جاءتكم البينات)}
ذكر \rawi{الطبري} أن الزلل هنا هو العدول عن الحق بعد ظهوره، والانحراف عن طاعة الله بعد مجيء الحجج البينات. وفي هذا تهديد شديد لمن عرف الحق ثم تركه، لا لمن خفي عليه الحق ابتداء.

\section{تأويل (هل ينظرون إلا أن يأتيهم الله في ظلل من الغمام)}
وقف \rawi{الطبري} عند هذا الوعيد، فبين أنه تهديد لمن أخر الدخول في الحق بعد وضوح البينات، حتى كأنه لا ينتظر إلا مجيء الأمر الفاصل والعذاب القاطع.

\section{تأويل (سل بني إسرائيل كم آتيناهم من آية بينة)}
في هذا الموضع سلّى الله نبيه صلى الله عليه وسلم بذكر حال من قبله من الأمم، وأن بني إسرائيل قد جاءتهم آيات كثيرة بينات ثم بدلوا وكفروا وخالفوا.

ونبّه الشيخ هنا إلى معنى عظيم: أن القرآن لا يذكر هذا التاريخ لمجرد الحكاية، بل ليواسي المؤمن والداعي، ويضع له تكذيب الناس وردود أفعالهم في موضعها الصحيح. فليس من جاء بالبينات ثم كُذب بدعاً في تاريخ الرسالات.

\begin{fawaid}[من أعظم مقاصد قصص المكذبين تسلية المؤمن وتثبيته]
إذا علم الداعي أن الرسل أنفسهم جاءوا بالآيات البينات ثم كذبهم أقوامهم، خفّ عنه استغراب رد الناس للحق، وانصرف جهده إلى ما كلف به: البيان والبلاغ والصبر، لا إلى التحسر على كل من أعرض.
\end{fawaid}

\section{تأويل (ومن يبدل نعمة الله من بعد ما جاءته)}
فسر \rawi{الطبري} نعمة الله هنا بالإسلام وما أنزله من شرائع الدين، وأن تبديلها يكون بالكفر بها أو بترك العمل بها بعد مجيئها ووضوحها.

\separator

\section{تأويل (زين للذين كفروا الحياة الدنيا ويسخرون من الذين آمنوا)}
بيّن \rawi{الطبري} أن الكفار زينت لهم الدنيا فركنوا إليها، وصارت مقاييسهم مبنية على الزينة والرياسة والجاه، لذلك كانوا يسخرون من أهل الإيمان الذين تركوا كثيراً من زخارف الدنيا طلباً لما عند الله.

\section{تأويل (والذين اتقوا فوقهم يوم القيامة)}
قرر \rawi{الطبري} أن العلو الحقيقي يوم القيامة لأهل التقوى، لا لأهل الدنيا وزخرفها. وأشار الشيخ إلى أن هذا من مواساة الله لأوليائه، وأن القرآن كثيراً ما ينتصر للمؤمنين الضعفاء أمام استعلاء أهل الباطل.

\begin{fawaid}[العلو الحقيقي لأهل التقوى لا لأهل الزينة]
من سنن القرآن أنه يرد موازين الناس الفاسدة إلى الميزان الحق؛ فليس الشرف في كثرة المال ولا في جاه الدنيا، بل في التقوى التي يرفع الله أهلها فوق المستكبرين يوم القيامة.
\end{fawaid}

\section{تأويل (والله يرزق من يشاء بغير حساب)}
ذكر \rawi{الطبري} أن هذا من تمام المدح لله تعالى؛ لأنه يعطي من يشاء من غير خوف من نفاد خزائنه، ولا حاجة إلى حساب يضبط به ما يخرج من ملكه مخافة النقص.

\separator

\section{تأويل (كان الناس أمة واحدة فبعث الله النبيين)}
\isnad{قال أبو جعفر:} اختلف أهل التأويل في معنى \textit{الأمة} هنا، وفي الوقت الذي كان الناس فيه أمة واحدة:
\begin{itemize}
    \item فقيل: كانوا على دين الحق من آدم إلى نوح ثم اختلفوا.
    \item وقيل: المراد آدم نفسه إماماً لذريته.
    \item وقيل: كانوا أمة واحدة حين أُخرجوا من ظهر آدم.
    \item وقيل: بل كانوا أمة واحدة على الكفر.
\end{itemize}

وبعد عرض هذه الأقوال رجح \rawi{الطبري} أنهم كانوا أمة واحدة على الإيمان والحق، واستدل لذلك بآية سورة يونس: \begin{ayah}وَمَا كَانَ النَّاسُ إِلَّا أُمَّةً وَاحِدَةً فَاخْتَلَفُوا\end{ayah}، فدل على أن المذموم هو الاختلاف الطارئ، لا الاجتماع الأول.

ومع ذلك شدد على أمر منهجي مهم: وهو أن تحديد الزمن الدقيق الذي وقع فيه هذا الاجتماع الأول لا تقوم عليه حجة قاطعة من كتاب ولا سنة ثابتة، فلا يجوز الجزم فيه بما لا دليل عليه.

\begin{fawaid}[ليس كل ما ذكر في التفسير يجوز الجزم به]
إذا لم يكن في المسألة نص ثابت أو دلالة قاطعة، فغاية ما يقال فيها هو الاحتمال أو النقل عن بعض السلف، لا الجزم والقطع. ومن دقة الطبري أنه يميز بين ما ترجح معناه، وبين ما لا يضر الجهل بتعيين وقته أو شخصه.
\end{fawaid}

\section{تأويل (وأنزل معهم الكتاب بالحق ليحكم بين الناس فيما اختلفوا فيه)}
بيّن \rawi{الطبري} أن إضافة الحكم إلى الكتاب إنما هي من جهة أن الأنبياء يحكمون بما دلهم عليه الكتاب المنزل، فكأن الكتاب هو الحاكم بين الناس بدلالته وإرشاده، وإن كان الفاصل في الخصومات هو الرسول أو النبي.

\separator

\section{تأويل (يسألونك عن الشهر الحرام قتال فيه)}
ذكر \rawi{الطبري} الخلاف في حكم القتال في الأشهر الحرم، ثم رجح أن النهي المتقدم منسوخ بما نزل بعده من آيات القتال العامة، واستدل على ذلك بالسيرة النبوية، وبوقوع غزوات النبي صلى الله عليه وسلم ومبايعته على القتال في بعض الأشهر الحرم.

ونبّه الشيخ هنا إلى فائدة مهمة: أن \rawi{الطبري} من المفسرين الذين يضيقون باب النسخ ولا يتوسعون فيه، فإذا صرح بالنسخ في موضع كان ذلك من المواضع الجديرة بالتقييد والنظر.

\begin{fawaid}[تضييق باب النسخ لا يمنع إثباته عند قيام دليله]
منهج الطبري في النسخ منهج متحفظ، ولذلك فإثباته النسخ في موضع من المواضع ليس أمراً عابراً، بل يكون مبنياً عنده على دليل من القرآن أو السنة أو السيرة الثابتة. وهذا يعلّم القارئ ألا يكثر دعوى النسخ بلا برهان.
\end{fawaid}

\section{تأويل (ولا يزالون يقاتلونكم حتى يردوكم عن دينكم إن استطاعوا)}
فسر \rawi{الطبري} الآية بأن المشركين لا يزالون قائمين على غايتهم في رد المؤمنين عن دينهم ما وجدوا إلى ذلك سبيلاً، فحذر الله عباده من عداوتهم، وكشف لهم حقيقة مقصدهم.

\section{تأويل (ومن يرتدد منكم عن دينه فيمت وهو كافر)}
بيّن \rawi{الطبري} أن الإحباط المذكور هنا إنما هو لمن مات على الكفر بعد الردة، أما من تاب ورجع قبل الموت فباب التوبة مفتوح له.

\separator

\section{تأويل (إن الذين آمنوا والذين هاجروا وجاهدوا في سبيل الله)}
ختم \rawi{الطبري} هذا المقطع بذكر الصنف المقابل لمن ارتد أو ضعف: أهل الإيمان والهجرة والجهاد، وهؤلاء هم أرجى الناس لرحمة الله.

ونبّه الشيخ إلى أن هذه الثلاثية ليست محصورة في صورتها التاريخية الخاصة بالصحابة فقط، بل لها معنى عام في كل مؤمن:
\begin{itemize}
    \item آمن بالله حقاً.
    \item وهاجر ما نهى الله عنه.
    \item وجاهد نفسه وعدو الله بقدر استطاعته.
\end{itemize}

\begin{fawaid}[الإيمان والهجرة والجهاد معانٍ باقية في الأمة]
ليست الهجرة والجهاد معاني منقطعة بانقطاع زمن معين، بل لكل مؤمن نصيب منهما بحسب حاله: يهجر ما نهى الله عنه، ويجاهد نفسه على الطاعة، ويثبت على الإيمان، فيدخل في هذا الوعد بقدر صدقه.
\end{fawaid}

\separator

ختم الشيخ اللقاء بالتأكيد على أن أعظم أصول التفسير وفهم القرآن ما أثر عن النبي صلى الله عليه وسلم وصحابته، وأن كتب السنن والآثار من أقوى ما يفتح لطالب العلم أبواب البيان في القرآن، كما أن القرآن نفسه دواء للشكوك، وتثبيت للمؤمن، وتصحيح لأولوياته في الدعوة والبلاغ والصبر.
